\section{Introduction}
\label{sec:intro}

\IEEEPARstart{L}{oRa} and LoRaWAN support long-range, low-cost connectivity for smart-city sensing, industrial monitoring, and other large-scale Internet-of-Things (IoT) deployments~\cite{hamdaoui2022deeplearning}.
Because many endpoints are resource-constrained, physical-layer radio frequency fingerprint identification (RFFI) offers a lightweight complement to cryptographic authentication at the gateway: a receiver can verify that over-the-air IQ captures match an enrolled transmitter identity from hardware-induced impairments~\cite{brik2008wireless_device_id,jian2020deeplearning_rf_fingerprinting,xie2021physical_layer_auth_survey}.

Cross-receiver deployment is a critical practical challenge for LoRa RFFI.
Transmitter fingerprints are not observed directly; they are filtered by the receiving front-end.
When the gateway radio changes, receiver-specific spectral response, gain imbalance, and filtering can distort the observed impairment pattern even if the transmitter set and environment remain unchanged~\cite{elmaghbub2021lora_wild,ma2025receiver_independent_rffi}.
As a result, models trained on one receiver often fail on another.
Prior work on the public OSU LoRa corpus reports that strict source-only cross-receiver transfer can remain near chance despite strong same-receiver performance on hybrid models from a related submitted manuscript~\cite{paper1_placeholder,elmaghbub2021lora_wild,shen2021lora_jsac}.

Out-of-band (OOB) spectral leakage has long been recognized as a source of transmitter hardware evidence in LoRa RFFI~\cite{elmaghbub2020widescan,shen2021lora_jsac,hamdaoui2022deeplearning}.
In a related submitted manuscript, OOB evidence is fused with in-band representations through cross-attention on a fixed receiver~\cite{paper1_placeholder}.
Our diagnosis, however, reveals a dual role: under a fixed receiver, OOB features can carry device-discriminative information, whereas under cross-receiver shift the OOB path becomes \emph{more receiver-entangled} than the main in-band path.
This receiver-induced OOB feature entanglement helps explain why source-only transfer remains poor even when a strong same-receiver backbone is available.

Motivated by this diagnosis, we study lightweight target-receiver calibration rather than redesigning the backbone.
We show that a source-trained classifier and a source-only prototype are not transferable under cross-receiver shift, whereas target-receiver local prototypes built from a small number of labeled calibration windows per device can restore transmitter separability.
Unsupervised target adaptation baselines---including entropy minimization and pseudo-label prototype updates on an unlabeled calibration pool---remain insufficient under severe cross-receiver collapse, even after confidence-threshold tuning.

This paper proposes a diagnosis-first framework and Receiver-Calibrated Prototype Adaptation (RCPA).
We use a frozen OOB-guided RF-HSTU hybrid model from a related submitted manuscript only as a feature extractor~\cite{paper1_placeholder}; \textbf{we do not claim a new backbone architecture or a published prior method}.
Our focus is cross-receiver failure analysis and post-hoc target-receiver calibration.

\subsection{Contributions}
The main contributions are:
\begin{enumerate}
  \item We diagnose receiver-induced OOB feature entanglement in LoRa cross-receiver RFFI using OOB spectral profiling, receiver/device linear probes, embedding-distance analysis, and confusion-collapse analysis.
  \item We show that source-trained classifier boundaries and source-only prototypes are not transferable across receivers, and propose RCPA-T, which constructs target-receiver class prototypes from $K$ labeled calibration windows per device.
  \item We validate RCPA under a block-disjoint calibration/support/query protocol across two transfer directions, three checkpoint seeds, and three split repeats, with auxiliary analyses on OOB representation equalization and unsupervised test-time adaptation baselines.
\end{enumerate}

\textbf{Scope note:} $K$ denotes the number of labeled calibration \emph{windows} per device, not independent capture files.
Each device contributes one capture file per receiver in our dataset; windows are sampled deterministically from that file.
All calibration, support, and query windows are block-disjoint; query windows are never used for prototype construction or adaptation-statistic estimation.

The remainder of this paper is organized as follows.
Section~\ref{sec:related} reviews related RFFI, cross-receiver, adaptation, and prototype-classification work.
Section~\ref{sec:diagnosis} presents the cross-receiver failure diagnosis.
Section~\ref{sec:rcpa} describes RCPA and the block-disjoint evaluation protocol.
Sections~\ref{sec:experiments}--\ref{sec:results} report experimental setup and results.
Sections~\ref{sec:discussion}--\ref{sec:conclusion} discuss limitations and conclude.
